%%%%%%%%
\section{Relevant aspects of the Gorelik--Kac determinant} \label{sec:GKvacdet}
%%%%%%%%
In this section, we review the Gorelik--Kac determinant of the universal vacuum module for simply-laced affine Kac--Moody algebras \cite{gorelik2007simplicity}. Following the simplified form presented in \cite{ArabiArdehali:2025fad}, the vacuum determinant for $\gf$ simply laced is given by, 
\begin{equation} \label{eq:KacGorelikvacdet}
    \text{det}~S_\nu(k) \propto (k+h^\vee)^{m_0(\nu)} \prod_{\substack{p \geqslant 2, q \geqslant 1 \\ (p,q) = 1}} (k-k_{p,q})^{m_{p/q}(\nu)}~,
\end{equation}
where $h^\vee$ is the dual Coxeter number of $\gf$ and the levels are defined by,
\begin{equation}
    k_{p,q} = - h^\vee + \frac{p}{q}~.
\end{equation}
The generating function of the exponent $m_0(\nu)$ is given by,
\begin{equation}
    \sum_{\nu \in \hat{\Delta}_+} m_0(\nu) e^{-\nu} = \frac{1}{\tilde{R}} \sum_{m = 2}^{h^\vee} \sum_{i = m}^\infty (s_{m-1} - s_m) \frac{e^{-i \delta}}{1 - e^{-i \delta}}~,
\end{equation}
where $\delta$ is the imaginary root of $\hat{\gf}$, $s_j := \# \{ \alpha \in \Delta_+ : (\rho|\alpha) = j \}$, $\rho$ is the Weyl vector of $\gf$, $\Delta_+$ is the set of positive roots of $\gf$, and
\begin{equation}
    \tilde{R} = \sum_{\gamma \in \hat{\Delta}_+ / \Delta_+} (1-e^{-\gamma})^{\text{dim} \hat{\gf}_\gamma}~.
\end{equation}
Here, $\hat{\Delta}_+ = \Delta_+ \sqcup \left\{\alpha + n\delta\,\middle|\,\alpha \in \Delta, \ n\geqslant 1 \right\} \sqcup \left\{n\delta \,\middle|\, n \geqslant 1 \right\}$ is the set of positive affine roots and dim$\hat{\gf}_\gamma$ = dim$\hf$ if $\gamma \in \mathbb{Z} \delta$, and dim$\hat{\gf}_\gamma$ = 1 if $\gamma \notin \mathbb{Z} \delta$. 

The generating function for the exponents $m_{p/q}(\nu)$ is given by,
\begin{equation} \label{eq:multAKMmain}
    \sum_{\nu \in \hat{\Delta}_+} m_{p/q}(\nu) e^{-\nu} = \frac{M_{p/q}}{R}~,
\end{equation}
where,
\begin{equation}
\begin{split}
    &M_{p/q} := \sum_{\alpha \in \Delta_+} \sum_{j \in \mathbb{Z}_{\geqslant 1}}^{\prime} e^{-(jp-(\rho|\alpha))jq \delta} E((jp - (\rho|\alpha))\alpha) - \sum_{\alpha \in \Delta_+} \sum_{j \in \mathbb{Z}_{\geqslant 1}} e^{-(jp+(\rho|\alpha))jq \delta} E(jp\alpha)~, \\ 
    &R := \prod_{\gamma \in \hat{\Delta}_+} (1-e^{-\gamma})^{\text{dim} \hat{\gf}_\gamma}~, \\ 
    &E(\lambda) := \sum_{w \in \text{Weyl}(\gf)} (-1)^{\ell(w)} e^{w(\lambda + \rho)-\rho}~.
\end{split}
\end{equation}
The prime on the first sum indicates that the sum is over all $j \in \mathbb{Z}_{\geqslant 1}$ satisfying $j p -(\rho|\alpha) \geqslant 1$. It was shown in \cite{gorelik2007simplicity} that the exponent $M_{p/q} = 0$ for $p = 1$ and $p < 0$. A key ingredient in computing the exponents $M_{p/q}$ explicitly is the following regularity criterion:
\begin{equation} \label{eq:ADEregularity}
    E(\lambda) \neq 0 \Longleftrightarrow \lambda+\rho \ \ \mathrm{is} \ \ \mathrm{regular}, \ \ i.e., \forall \ \alpha \in \Delta_+, \ (\lambda + \rho|\alpha) \neq 0~. 
\end{equation}
To keep track of the different contributions to $M_{p/q}$, it is convenient to decompose it as $M_{p/q} = M_{p/q}^+ - M_{p/q}^-$ where,
\begin{equation} \label{eq:splitM}
\begin{split}
    &M_{p/q}^+ = \sum_{\alpha \in \Delta_+} \sum_{\substack{t \in \mathbb{Z}_{\geqslant 1}: \\ p|(t+(\rho|\alpha))}} e^{-t(t+(\rho|\alpha))\frac{q} {p}\delta} E(t \alpha) =: \sum_{\nu \in \hat{\Delta}_+} a^+_\nu e^{-\nu}~, \\
    &M_{p/q}^- = \sum_{\alpha \in \Delta_+} \sum_{\substack{t \in \mathbb{Z}_{\geqslant 2}: \\ p|t}} e^{-t(t+(\rho|\alpha))\frac{q}{p}\delta} E(t \alpha)  =: \sum_{\nu \in \hat{\Delta}_+} a^-_\nu e^{-\nu}~.
\end{split}
\end{equation}
As with the notation used in \cite{ArabiArdehali:2025fad}, we denote the restriction of the exponent at a fixed conformal level $L$ by $M_{p/q,L} := M^{+}_{p/q,L} - M^{-}_{p/q,L}$ where $M^\pm_{p/q, L} = \sum_{\mu \in Q^+} a^\pm_{L\delta + \mu} e^{-\mu}$, with $Q^+$ denoting the positive root lattice of $\gf$.  

The computation of the multiplicities $M_{p/q,L}$ is facilitated by the determination of the various $(t,\alpha)$ appearing in \eqref{eq:splitM} with $E(t\alpha) \neq 0$. In particular, note that they must satisfy $B_{(t,\alpha)} := t(t+(\rho|\alpha)) \frac{q}{p} \in \mathbb{N}q$ and $B_{(t,\alpha)} \leqslant L$. 

The strategy used in determining the multiplicities of zeros in the vacuum determinant will be to first characterise the possible values of $A_{(t,\alpha)} := t(t+(\rho|\alpha))$ for which $E(t\alpha)\neq0$, note down the corresponding $(t,\alpha)$, and then use that to compute relevant multiplicities. 

The multiplicities of the zeros of the Gorelik--Kac vacuum determinant (in all charge sectors) for simply-laced affine Kac--Moody algebras are encoded in
\begin{equation}
    \frac{M_{p/q,L}}{E(0)} = \frac{M^+_{p/q,L}}{E(0)} - \frac{M^-_{p/q,L}}{E(0)}~,
\end{equation}
where $E(0) = \prod_{\alpha \in \Delta_+}(1-e^{-\alpha})$. In general, the regularity criterion \eqref{eq:ADEregularity} and the expression for the exponents in \eqref{eq:splitM} imply that these functions are expressible as sums of characters of representations of the respective Lie algebra $\gf$,\footnote{While each weight $\lambda = t\alpha$ appearing in the expressions \eqref{eq:splitM} is integral, it need \emph{not} be dominant. However, the regularity criterion in \eqref{eq:ADEregularity} implies that one can perform a series of Weyl reflections so that the weight obtained after the shifted Weyl action $w.\lambda = \Lambda$ is dominant, and as a consequence $E(\lambda) = (-1)^{\ell(w)} E(\Lambda)$.} 
\begin{equation} \label{eq:expinchar}
\begin{split}
    \frac{M_{p/q,L}}{E(0)} &= \sum_{\substack{\Lambda \ \mathrm{dominant} \ \mathrm{integral} }} s_\Lambda \frac{E(\Lambda)}{E(0)}\\ &=\sum_{\substack{\Lambda \ \mathrm{dominant} \ \mathrm{integral} }} s_\Lambda~\mathrm{ch}(\Lambda), \qquad s_\Lambda \in \mathbb{Z}~. 
\end{split}
\end{equation}
Here $\mathrm{ch}(\Lambda)$ is the character of the representation of $\gf$ with highest weight $\Lambda$. Therefore, determining the multiplicities of zeros in the Gorelik--Kac vacuum determinant in a fixed-weight sector at a fixed conformal level reduces to finding the $(\Lambda, s_\Lambda)$ appearing in the above function. As a simple example, note that for $k = k_{p,q}$ boundary admissible, \emph{i.e.}, $p = h^\vee$ and $(p,q) = 1$,
\begin{equation}
    \frac{M_{p/q,q}}{E(0)} = \frac{E(\theta)}{E(0)} = \mathrm{ch}~\mathrm{Ad}~,
\end{equation}
which implies that the conformal level $L = q$ primitive null state at $k = k_{p,q}$ is a component of the adjoint representation. 